\vspace{0.8em}
\section{国内外研究现状及分析}\label{Sec: literature review}
\vspace{0.8em}

螺旋桨的气动外形决定其飞行性能，从而决定了无人机的飞行效率。气动外形的优化方法是设计螺旋桨外形的重要研究方向。优化方法的流程可分为五个阶段：（1）基准模型的选择；（2）目标函数以及优化参数的确定；（3）代理模型的建立；（4）在约束条件下进行寻优；（5）优化结果的验证。其中确定目标函数以及优化参数、建立代理模型和寻优过程是整个优化设计过程中的难点和关键。
目标函数和优化参数决定整个优化过程的方向以及最终地优化结果。螺旋桨气动外形优化的目标函数一般为最小化阻力系数、最大化升力系数或最大化升阻比。
确定目标函数和优化参数涉及到如何对螺旋桨及其旋翼进行参数化研究。
代理模型是用来在优化过程中计算螺旋桨的空气动力学参数的模型，
代理模型的准确与否关系到优化数据的精度、计算效率。建立代理模型的方法主要分为数值计算模型和基于人工智能的代理模型。
寻优过程涉及到如何选择优化方法，如何分析优化数据，如何确保数据具有较高精度等问题。

基于上述分析，本节将从以下主要介绍螺旋桨气动外形优化方法的研究，主要分为几个方面：

（1）螺旋桨外形参数化方法研究。

（2）螺旋桨气动参数代理模型建立方法研究。

（3）多重保真度模型的研究。

（4）螺旋桨气动外形优化方法的研究

（5）基于拓扑特征的优化数据分析研究。


\vspace{0.8em}
\subsection{国内外螺旋桨的外形参数化方法研究现状}\label{Sec: shape parameter}
\vspace{0.8em}

螺旋桨的外形参数主要包括叶片数目、叶片直径、翼型、叶片扭转分布、桨毂直径等。小型四旋翼无人机用的螺旋桨的叶片数目主要为两个或三个，叶片直径一般为几十厘米，桨毂直径一般为几厘米。叶片数目、叶片直径、桨毂直径一般在基准模型确定时已为定值。用于优化的输入参数一般为叶片的扭转分布或描述螺旋桨翼型的几何参数。叶片的扭转分布即为叶片不同位置扭转角的大小。翼型的外形由两条曲线组成，如何描述这两条曲线则是参数化方法的研究重点。

Masters等人\cite{Dominic2015, Masters2017}、Sripawadkul等人\cite{Vis2010}和Salunke等人\cite{Salunke2014}对翼型的参数化方法做了总结和对比分析。翼型的参数化方法主要有Ferguson曲线方法\cite{Ferguson1964}、Hicks-Henne型函数方法\cite{Hicks1978}、B样条曲线方法\cite{Farin1997}、特征参数描述（PARSEC）方法\cite{Sobieczky1999}、自由曲面变形方法（Free Form Deformation, FFD）\cite{Thomas1986}、类别形状变换方法（Class function/Shape function Transformation, CST）\cite{Kulfan2006}。

Ferguson曲线方法需要6个参数来构建翼型的形状曲线，该方法简单\cite{Sun2018}，但是由于使用的参数较少，所以精度不高。
Hicks-Henne型函数方法通过在基准翼型上叠加扰动型函数来描述新的翼型的形状\cite{Zhang2011}。
B样条曲线方法是一种广泛使用的生成分段多项式曲线的方法，在基于CAD外形优化的设计者中广泛使用\cite{Xu2014, Martn2014}。
PARSEC方法采用11个几何特征参数去描述翼型的几何形状，由于其利用的特征参数较少，对后缘位置，即气动特性影响较大位置的形状参数化近似效果不好\cite{Salunke2014}，因此精度较低\cite{Shahrokhi2007, Kharal2012}。针对后缘形状的参数化，Sobieczky\cite{Shahrokhi2007}于2007年给出了一定修正，使后缘形状近似于凹包。
FFD通过将变形物体嵌入一个简单柔软的实体中，实体变形带动嵌入物发生形变\cite{Thomas2012,Leloudas2018, Hu2022Shape}。
Kulfan提出的CST方法包括类别函数和形状函数两部分，类别函数控制气动外形的一般形状，形状函数用来增加扰动改变翼型形状。因其能够使用较少的设计变量即可进行翼型参数化，使用及其灵活，因而应用广泛\cite{Wu2022,Tan2022,Zhang2021}。

近年来，部分研究人员通过结合多种翼型参数化方法，从而实现减小误差。如Hu等人将B样条方法和CST方法相结合，以较少的参数实现翼型的几何优化\cite{Hu2022}。邬晓静等人将CST方法和POD方法结合，同样实现了以较少参数进行翼型气动外形优化的目的\cite{Wu2019}。

\vspace{0.8em}
\subsection{国内外螺旋桨的气动参数代理模型建立方法研究现状}\label{Sec: response model}
\vspace{0.8em}

在优化问题过程中建立代理模型，本质上是探索问题空间的某种内部模型。目前，基于代理模型的优化（Surrogate-based optimization, SBO）\cite{Forrester2009,Queipo2005}是目前螺旋桨外形优化的最有效方法之一。
代理模型用于在优化过程中计算新气动外形的目标函数，它是建立在输入变量与输出量之间的分析模型，也称为近似模型、响应模型。
代理模型能够代替复杂耗时的实验或CFD计算，从而大幅缩减优化过程中计算新优化模型气动参数的时间，在优化问题中应用广泛。
建立一个准确且快速的代理模型十分重要。如果代理模型的结果误差非常大，那么它不能代替实验或高仿真的CFD计算得到新几何外形的空气动力学参数结果。在整个优化过程中，一般要经过至少近百次的外形迭代，如果代理模型的计算较慢，那么整个优化周期就会大幅增长。
与此同时，代理模型的试验设计方法\cite{Queipo2005}、代理模型的推广能力\cite{Fornberg2004,Rippa1999}、代理模型的评价方法\cite{Romero2005}等是在建立代理模型时关注的重点。

建立代理模型的基本步骤是（1）选取待优化的变量；（2）设计采样方法；（3）选取代理模型的种类并建立模型；（4）代理模型的验证。
在螺旋桨的外形优化问题中，待优化的变量通常为螺旋桨的几何外形，如叶片的扭转角\cite{Qiao2013, WangTY2023}等。
为保证采样的随机性，一般采用随机采样方法，如蒙特卡洛法（Monte Carlo）、拉丁超立方采样法（Lating hypercube sampling, LHS）~\cite{Iuliano2018,Liu2023}。
代理模型的种类可分为基于数值计算的方法和基于机器学习的方法。为了缩短计算时间，基于数值计算的方法一般采用低保真度的数值方法，通过编写代码或使用相关软件进行计算。基于机器学习的方法一般采用高保真度的数据训练一个快速响应、较为准确的模型，从而缩短在优化过程中的计算时间。
最后，代理模型的验证一般是通过代理模型的结果与实验数据的对比、或与高保真度的CFD仿真数据对比，采用如均方根误差（Root mean square error, RMSE）、决定系数$R^2$等的统计量进行准确性的评估。
建立代理模型的方法是关于代理模型研究的关键，下面主要探讨国内外研究人员建立代理模型的方法。

\vspace{0.8em}
\subsubsection{基于数值计算的代理模型建立研究现状}\label{Sec: numeric model}
\vspace{0.8em}

数值计算方法是通过数值方法近似求解微分、积分方程或方程组的方法，通常采用计算机程序进行计算求解。在螺旋桨的外形优化问题中，数值计算方法即为离散描述螺旋桨附近流场的空气动力学方程组，求解其相应的空气动力学参数。为保证代理模型能够快速得到新的计算结果，所以一般采用较为简单的、精度不太高的数值计算方法，主要的方法有升力线理论（lifting line theory）、
涡格法（Vortex Lattice Method，VLM）、
叶素动量理论（Blade Element Momentum Theory, BEMT）、基于有限体积法的CFD方法等。

升力线理论由普朗特于1913至1918年提出，该理论用一条附着于机翼上的集中涡线和从机翼顺流延伸到无穷远的尾涡流面来代替机翼作用的一种理论。该方法主要用于计算亚声速流动中大展弦比的升力和由尾涡引起的阻力。升力线理论由于计算简单快速，在机翼几何设计中用一定的应用。乔宇航等人从升力线理论推导出机翼几何扭转的设计方法，并通过CFD结果验证了该方法的准确性\cite{Qiao2013}。近年来，为提升升力线理论的计算精确度，一些研究人员对升力线理论进行修正或改进。比如，乔诗展等人采用深度学习方法进行改进\cite{Qiao2022}，Goates等人通过定义更有效的升力线使该方法能够在扁几何条件下得到较为准确的结果\cite{Goates2022}。

涡格法的首次出现时在1937年Pistolesi的报告中\cite{Pistolesi1937}，在经过几十年的发展，涡格法逐渐完善，从稳态计算到非稳态计算、非线性计算，准确性逐渐提升。涡格法在超音速飞行的计算和验证是由Miranda在1976年完成的\cite{Miranda1976}。美国国家航空航天局（National Aeronautics and Space Administration, NASA）、美国麻省理工学院（Massachusetts Institute of Technology,MIT）等曾基于涡格法开发相应的计算程序，如VORLAX\cite{Miranda1976}、AVL\cite{Budziak2015}、ABSOLUT\cite{Bueso2011}等。如今，涡格法已被广泛应用于螺旋桨的外形优化问题\cite{Cho1998,Kier2005,Guimares2021}。

叶素动量理论是叶素理论与动量理论的结合，通过引入诱导系数，迭代求解，计算出作用在螺旋桨力和力矩。1925年C.N.H. lock等人将叶素理论与动量理论结合使用，1926年由Hermann Glauert将该理论推导写成现代的形式。
叶素动量理论在发展过程中有大量的修正优化，如三维效应修正、叶尖损失修正、轮毂损失修正等。
叶素动量理论因计算快效率高而被广泛应用，如Matthew等人\cite{McCrink2017}用该理论预测固定螺旋桨性能，Umberto等人\cite{Boatto2023}比较了多种该理论的修正的计算结果，认为Zhong等人\cite{Zhong2020}提出的叶尖损失适用于文中使用的BREL 5 MW涡轮机。
由于叶素动量理论模型简单、计算效率高、适用性广、灵活性高，有大量的研究人员基于该理论开发了螺旋桨设计计算的应用软件和开源代码。如QBlade\cite{Marten2013QBLADEA}、JBlade\cite{Silvestre2013}、pyBEMT\cite{Giljarhus2020pyBEMT}、AeroDyn等。其中，QBlade的应用最为广泛\cite{Turbak2023, Vyom2022, Fernandez2023,Panjwani2020}。

基于有限体积法的CFD方法由于需要计算复杂的非定常的流体流动问题，往往计算较为缓慢。
因此，一些研究人员通过自己编写开发CFD计算程序并小范围地进行应用\cite{Wu2022, Shahrokhi2007}。
然而，由密歇根大学MDO实验室开发开源软件ADflow~\cite{Mader2020}被广泛用于飞机设计优化中\cite{Chauhan2021,Scco2021,Li2021}，因为该CFD软件采用欧拉方程计算可压缩流问题，离散方法采用结构化多区块
%（structured multi-block）
方法。

\vspace{0.8em}
\subsubsection{采用机器学习建立代理模型的研究现状}\label{Sec: surrogate model}
\vspace{0.8em}

数值计算方法为节省计算时间，往往采用低保真度的数值方法，其结果可以保证与实验数据或高保真度的CFD仿真数据有相同的趋势，但是误差较大。
随着计算机、机器学习理论和大数据理论的发展，为了进一步提高代理模型的准确性，基于机器学习方法训练得到的代理模型的方法被广泛应用。该方法需要用到大量的数据去训练模型，这些数据的来源一般为实验数据或高精度的CFD仿真计算数据。
训练代理模型所需的样本数以及训练时间、代理模型的泛化能力是基于机器学习理论建立代理模型的研究关键。

目前，研究人员尝试采用多种多样的机器学习方法来生成代理模型，这些方法主要有多项式响应面\cite{Ahn2001, Xiong2006}、移动最小二乘插值（Moving least-squares, MLS）\cite{Toropov2005}、径向基函数（Radial basis functions，RBF）\cite{Asouti2023}、人工神经网络（Artificial Neural Network, ANN）\cite{Liu2023,Jahangirian2011,Han2017, Zan2022}、Kriging模型\cite{Jouhaud2007,LiJichao2020,Wu2022}、支持向量机（Support vector regression, SVR）\cite{Andres2012}、强化学习（Reinforcement Learning, RL）\cite{Chen2023}、本征正交分解（Proper Orthogonal Decomposition, POD）\cite{Iuliano2013,Grey2018,LiJichao2019}等。

西北工业大学韩忠华教授开发一个名为“SurroOpt”的基于代理模型的空气动力学优化代码\cite{Han2016}，该代码中的代理模型包括应用广泛的机器学习方法，如Kringing、RBF、SVR、ANN等，该代码在小范围内得到应用\cite{Han2020,WangKai2023}。

一些研究人员采用联合多种方法的方式建立代理模型，从而在提高代理模型的精度的同时减少优化算法的计算量。比如，Wang Xu等人用RBF和神经网络结合生成代理模型，用来分析机翼空气动力学参数的不确定性\cite{Wang2019}。

%总结多少参数，精度多少，用什么数据。

\vspace{0.8em}
\subsection{国内外多重保真度模型的研究现状}\label{Sec: multi-fidelity}
\vspace{0.8em}

在上文分析的建立代理模型的方法中，为了保证优化过程耗时尽可能少，基于数值计算方法建立的代理模型牺牲了一些精确度，这是一种低保真度的模型。基于机器学习方法建立的代理模型因为用于训练生成模型的数据往往是精确度高的实验数据或高精确度的CFD仿真实验数据，前期获取数据的困难大、时间长，这是一种高保真度的模型。
现实中的螺旋桨外形优化问题面临着几何结构参数多、螺旋桨周围流场复杂，所以采取实验研究获取数据的方法耗费的人力物力财力多，如果CFD仿真计算，一个算例的计算时间往往以周为单位，高保真模型的建立周期长。多重保真度模型利用大量低保真度的数据来获取系统参数的基本趋势，再通过少量高保真度的数据提高结果的准确性，它既有低保真都模型计算快的优点，又有高保真度模型的精度高的优点，因而在近年来的螺旋桨气动外形优化问题研究中大受欢迎。
根据NASA的技术报告，多重保真度模型的开发需要长期研究，在2030年之后仍可能是一个活跃的研究课题\cite{Slotnick2014,Slotnick2014Enab}。

多重保真度模型依据数据融合的方式主要分为三类：（1）基于缩放函数的多重保真度模型；（2）空间映射的多重保真度模型；（3）Cokriging方法\cite{Zhou2023}。多重保真度模型根据应用可分为分层多重保真度模型和不分层多重保真度模型\cite{Feldstein2020}。近十年来，在航空航天领域内已经涌现出讨论如何建立并应用多重保真度模型帮助机翼或螺旋桨设计和优化的研究，建立多重保真度模型的方法主要有空间映射和Cokriging方法。

空间映射建立低保真度模型参数和高保真度的模型参数之间的关系\cite{Bakr2001}。Usama Toman等人采用该方法建立多重保真度的代理模型，用于飞机螺旋桨几何设计和优化问题中，结果表明该方法比用kringing建立代理模型要高效\cite{Toman2019}。

相比与空间映射，cokriging方法的应用较为广泛。Zhang Yu等人采用多级分层Kriging模型作为代理模型做RAE2822翼型和ONERA M6机翼得阻力最小化问题，结果表明采用多级分层Kriging模型进行优化分析能够更快收敛到与高保真度模型几乎相同的最小值\cite{Zhang2024}。Chenzhou Xu等人对比Cokriging模型与不分层Cokriging模型、线性回归多保真度模型、潜映射高斯过程模型进行对比，结果表明Cokriging模型高效准确\cite{ChenzhouXu2024}。

\vspace{0.8em}
\subsection{国内外螺旋桨的气动外形优化方法研究现状}\label{Sec: optimization}
\vspace{0.8em}

优化方法可分为解析方法和数值方法。解析方法通过严格的数学推导证明给出优化问题的解析解，但复杂的实际问题不能满足很多苛刻条件，因此难以实现。数值方法的发展得力于计算机的快速发展，它通过迭代不断逼近最优解，以得到解析最优解的近似解，在实际应用中更具有普遍意义。应用较为广泛的数值优化方法大多需要用到梯度信息，主要包括最速下降法、拟牛顿法、序列二次规划和共轭梯度法。这些方法的特点是速度快、局部搜索能力强，但是常会陷入局部最优解、适应性差。
随着人工智能、机器学习的发展，一些新型智能优化算法出现，它们大多受到生物进化、物理原理或群体行为的启发，不需要计算梯度。它们全局搜索能力强，适用于处理多目标优化问题，如遗传算法、模拟退火算法、粒子群算法等。
同时，在多年的发展与积累，在气动设计领域已积累的大量的螺旋桨翼型相关的几何、气动性能参数，并已形成如伊利诺伊大学香槟分校的Michael Selig教授主导整理的UIUC数据库，密歇根大学MDO实验室整理的Webfoil数据库等。数据库的涌现使得基于数据的人工智能算法在气动外形优化领域得以发展壮大\cite{Oliveira2023}。

螺旋桨气动外形优化问题中分为单目标优化问题和多目标优化问题。单目标优化问题的优化目标主要有最小化阻力系数、最大化升力系数、最大化升阻比、最大化$FM$（figure of merit）。
在单目标优化问题中常采用的优化方法有遗传算法、粒子群算法、深度强化学习、基于“教”与“学”的优化算法等。遗传算法是一种应用比较广且开始时间比较早的一种优化算法，SHAHROKHI等人通过遗传算法对分别对两种翼型进行最大化升阻比研究\cite{Shahrokhi2007}。
Jing Zhang等人通过遗传算法进行最大化翼型升阻比研究\cite{Zhang2021}，文中优化的翼型最大化升阻比为18.2543。Jiaqi Liu等人用遗传算法对SC1095翼型进行优化以应对失速抑制效应\cite{Liu2023}。
在粒子群算法的应用方面，刘坤澎等人采用粒子群算法以某太阳能无人机为例设计一种能够满足多设计点螺旋桨气动外形优化和多工况变桨距的方法以提高无人机的效率\cite{LiuKP2023}。Qingdu Tan等人用粒子群算法获取了参数化翼型时所需的最优Bernstein系数\cite{Tan2022}。
Ziyang Liu等人用深度强化学习方法最大化翼型的升力系数，解决了超临界机翼优化设计问题\cite{Liu2024}。
西北工业大学邬晓敬提出一种基于“教”与“学”的优化算法\cite{Wu2022}，用于最小化翼型的阻力系数，值得一提的是，该优化算法能够解决高达200维的优化问题。

多目标优化问题的优化目标通常是上述单目标优化目标中的两个及以上的组合。目标函数之间是矛盾的，因而满足所有优化目标的解不唯一，而是一组解，被称为Pareto前缘。
遗传算法、粒子群算法等也可用于多目标优化问题，如Oliveira等人用遗传算法进行两种不同螺旋桨优化策略\cite{Oliveira2023}，一种为最大化升力系数的同时最小化弯曲系数，另一种为同时最大化升力系数、最小化弯曲系数、最小化螺旋桨体积和最小化效率。最终Oliveira等人总结了不同的优化策略所得到的不同的Pareto前缘，供决策者参考。Kitae Park等人也采用多目标的遗传算法进行多目标螺旋桨形状优化问题研究，新设计的螺旋桨做到了在降低能量消耗的同时提升效率\cite{Park2022}。
ShiYi Jin等人用粒子群算法探讨多种优化翼型的方法\cite{Jin2024}，作者分别进行了三种不同的优化目标，分别为最大化升力系数的同时最小化阻力系数、横向电场和横向磁场，最终第三种优化目标能够得到最大的升力系数。
Ye Jianchuan等人通过最大化无人机前进时间来优化无人机在前进状态下的表现，体现在优化螺旋桨翼型的弦长和扭转角分布，也就是改变螺旋桨的外形\cite{Ye2021}。

\vspace{0.8em}
\subsection{国内外优化数据分析研究现状}\label{Sec: data analysis}
\vspace{0.8em}

数据分析是通过一系列方法从数据中提取有价值的信息的过程。在螺旋桨外形优化过程中，有大量的螺旋桨周围流体流动的数据以及优化过程中产生的大量数据，如果能够从这些数据中提取到帮助加速优化过程的信息或最优解的相关信息并加以应用，那么优化周期将大幅缩短，优化的效果也可能会增强。

在螺旋桨外形优化过程中，为提取螺旋桨周围流场的数据特征，可以采用本征正交分解（Proper Orthogonal Decomposition，POD），动力学模态分解（Dynamic Modal Decomposition，DMD）等方法。
POD能够抓取流场数据中的主要模态，以低时间成本准确得还原流场\cite{Masters2017}，帮助在优化过程中筛选出气动性能好的几何，去除掉气动性能明显较差的几何\cite{Taol2010}。研究人员用POD提取的模态主要包括螺旋桨翼型的弯度\cite{Poole2015,Li2019}。因而，POD作为一种高效的螺旋桨气动外形设计的分析工具被广泛应用。Berguin和Mavris\cite{Berguin2014, Berguin2015}应用POD提取低维流动数据的模态，并发现该模态对空气动力学参数有显著影响。该方法能够帮助进行数据降价，并进行以增大升力或减小阻力的优化研究\cite{Grey2018,LiJichao2019}。
DMD是一种从非定常流场中提取流场中主要特征的数据驱动算法， 它能同时得到模态特征和动力学信息，因而具有时空耦合的优势。
Mariappan\cite{Mariappan2014}和Mohan\cite{Mohan2015}等人通过对比POD和DMD方法在二维翼型和三维机翼的动失速现象，发现DMD模态得到的流动结构在描述流动不稳定模态具有一定优势。寇家庆\cite{KouJQ2016}等人分别用POD和DMD方法研究翼型跨声速抖振模态分析，结果表明在激波间断处附近的流场使用POD效果较好，但是在非定常流场的模态提取问题中DMD更精确。

然而，当优化数据只有空气动力学参数数据，而非完整的流场数据时，如何分析待优化的数据以得到有利的信息是十分困难的。这时，数据拓扑结构的分析是一种有利的手段。数据的拓扑结构可以揭示数据形状，确定数据中极值点的数量和位置，尤其是高维数据的极值点信息，这些信息能够帮助选择合适的优化算法。这一领域的研究基础是离散场拓扑分析\cite{Edelsbrunner2002}，并已在流动分析中得以应用\cite{Kasten2016am}。离散场拓扑分析的基础是原始离散数据，该方法能够在不使用滤波等方法的情况下过滤掉不真实的极值点等拓扑特征，保证在不破坏原始数据的拓扑结构的条件下，提取数据中的拓扑特征，从而指导优化算法的选择。目前，将分析优化数据的拓扑结构的方法用于螺旋桨气动外形优化问题的研究是空白状态。

%\vspace{0.8em}
\subsection{国内外文献综述的简析}\label{Sec: status analysis}
\vspace{0.8em}

综合上述国内外研究现状来看，螺旋桨的气动外形优化设计研究随着螺旋桨的应用场景不断变化，相关研究开始得较早，也较为深入。随着大数据时代的到来，研究人员积极将机器学习、数据驱动等方法融入到螺旋桨外形优化得研究中去。在这个过程中，产生大量的数据，也需要处理大量的数据，虽然最终的优化结果非常好，但是优化效率并不是很高。
气动外形优化设计效率的影响因素主要有三个：（1）设计几何空间的维数；（2）空气动力学分析的耗时；（3）优化算法的收敛性。先进的人工智能算法能够帮助提高气动外形参数化的效率，减少计算时间的同时准确预测气动性能参数。下面具体介绍该方法的主要现状。

螺旋桨外形的几何参数个数可以是无限多的，这就导致优化过程的数据维数相当高，这对优化过程中用到的代理模型的构建方法、优化算法的寻优能力、计算机的数据计算能力提出相当高的挑战。同时在优化的过程中，优化算法会不断产生新的几何外形参数，随着优化参数的增多，一些形状异常的几何参数数量也会增加，比如波浪形转的翼型表面。这些异常形状的几何参数最终会导致优化算法无法找到最优解。目前已知的比较成功的高维优化算法是基于“教”与“学”的优化算法\cite{Wu2022}，能够解决高达200维的翼型优化问题。但该方法是否能更广泛地应用于其他优化问题还有待研究。设计几何空间的维数虽然可以无限大，但是如果能确定一个足够大的维数，即可较好的完成螺旋桨几何外形的优化任务，而不是无限制的增大数据维数，可能是一个比较合理的方法。
并且，解决高维问题需要巨大的计算资源，如果能够有把高维问题分解为几个更低维数问题的叠加则更有利于优化问题的解决。

基于CFD的优化计算需要计算复杂的流场信息，有时还需要计算非定常螺旋桨转动的周围流场数据，计算量巨大，计算一个算例的时间可能以天为单位。如果数据中有几个算例计算没有收敛或存在其他问题，则导致优化算法在数据集空间中无法找到最优解。更为高效快速的计算方法亟待开发，总结归类公开当前已有的数据也能够帮助加速采用机器学习算法解决螺旋桨优化问题。多重保真度模型既有低保真模型的计算快的优点又有高保真模型精确的优点，是一种能够快速准确获取数据的方法，目前在螺旋桨外形优化问题中主要采用Kriging方法进行多重保真度模型的建立\cite{Zhang2024, ChenzhouXu2024}，但是该方法忽略了低保真度模型与高保真都模型之间的几何关系，在高维问题的处理上效果不显著。

优化算法是否能够收敛取决于计算域数据的质量、优化算法本身的探索和挖掘能力。目前确保数据质量的方法主要为统计学方法，比如计算CFD仿真结果与实验数据结果的相对误差。统计学方法直观、操作性强，但是难以挖掘数据中隐藏的信息。同时，在采用机器学习算法时，同一组数据集采用不同的算法构建代理模型都能够得到非常好的拟合结果，难以选出最优的方法。这时分析数据的拓扑特征能够给出统计学方法不能得到的数据信息，帮助分析数据质量，同时给出选择优化算法的建议。
